\documentclass[11pt]{article}
\usepackage[a4paper,margin=2.4cm]{geometry}
\usepackage{graphicx}
\usepackage{amsmath}
\usepackage{amssymb}
\usepackage{booktabs}
\usepackage{caption}
\usepackage[hidelinks]{hyperref}
\usepackage{xcolor}
\usepackage{titlesec}
\usepackage{authblk}
\captionsetup{font=small,labelfont=bf}
\titleformat{\section}{\normalfont\large\bfseries\color{black!85}}{\thesection}{0.6em}{}
\titleformat{\subsection}{\normalfont\normalsize\bfseries\color{black!75}}{\thesubsection}{0.5em}{}
\setlength{\parskip}{0.5em}
\setlength{\parindent}{0pt}
\renewcommand{\baselinestretch}{1.05}

\title{\textbf{Climate-driven dengue early warning in Indonesia:\\
an East Java case study with a seven-province comparison}}
\author[ ]{Serius Miliyani Dwi Putri}
\affil[ ]{\small M.Ked.Trop (Epidemiology, Tropical Medicine)\\
Research Manager, Climate and Health Centre, Universitas Brawijaya\\
\texttt{Proof-of-concept analysis}}
\date{\small 2025}

\begin{document}
\maketitle

\vspace{-1.2em}
\noindent\rule{\textwidth}{0.4pt}
\begin{abstract}
\noindent Dengue transmission in Indonesia is strongly climate-sensitive, yet
routine surveillance remains largely reactive. As a compact, reproducible
demonstration of the modelling that underpins climate-based early-warning
systems, I link open monthly dengue surveillance (OpenDengue v1.3) with open
climate reanalysis (Open-Meteo / ERA5) and fit negative-binomial distributed-lag
non-linear models (DLNMs) of dengue \emph{incidence}. A single-province analysis
for East Java (2004--2024) is extended to a pooled seven-province model spanning
Java and Bali, with a population offset and province-specific baselines. Across
all provinces the strongest rainfall signal occurs at a delay of two to three
months, consistent with the biological lag between rainfall, vector breeding and
onward transmission. The pipeline is intended as a portfolio artefact and the
foundation for a doctoral programme on operational dengue early warning.
\end{abstract}
\noindent\rule{\textwidth}{0.4pt}

\section{Background}
Dengue is among the fastest-growing vector-borne diseases worldwide, and
Indonesia carries one of the highest burdens in Southeast Asia. Temperature and
rainfall govern \emph{Aedes} survival, biting rates and the extrinsic incubation
period, with effects that unfold over weeks to months. That delayed structure is
precisely what a distributed-lag model quantifies, and it is what makes
forecasting with useful lead time possible. This note is a deliberately minimal
first step: real surveillance data, an open climate source, and a standard,
defensible model, released reproducibly.

\section{Data}
Monthly, province-level (Admin1) reported dengue case counts were extracted from
the OpenDengue database v1.3, which standardises publicly available surveillance
data. Seven provinces with near-complete monthly coverage were retained (East
Java, West Java, Central Java, DKI Jakarta, Banten, Bali and DI Yogyakarta),
spanning January 2004 to July 2024. Provincial population (BPS Statistics
Indonesia census and intercensal projections) was interpolated to monthly values
and used as a model offset, converting counts to incidence. Monthly mean
temperature, total rainfall and relative humidity were obtained from the
Open-Meteo Historical Weather API (ERA5-based reanalysis) at representative
provincial coordinates.\footnote{The figures shown here use a clearly-labelled
synthetic climate series so the pipeline runs without network access; replacing
it with real Open-Meteo data is a single scripted step. Figures generated from
synthetic climate are tagged accordingly.}

\section{Methods}
Case counts $Y_{p,t}$ in province $p$ and month $t$ were modelled with a
negative-binomial GLM,
\begin{equation*}
\log \mathbb{E}[Y_{p,t}] = \log N_{p,t} + \alpha_p
+ \mathrm{cb}_{\text{temp}}(x_{p,t},\ell)
+ \mathrm{cb}_{\text{rain}}(r_{p,t},\ell)
+ s(\text{month}_t) + f_p(t),
\end{equation*}
where $N_{p,t}$ is population (offset), $\alpha_p$ a province-specific intercept,
and each \emph{cross-basis} $\mathrm{cb}(\cdot)$ is a tensor product of a
natural-spline exposure--response basis and a natural-spline lag--response basis
over lags $\ell = 0,\dots,6$ months (Gasparrini 2011). Seasonality $s(\cdot)$
uses cyclic terms and the trend $f_p(t)$ a natural spline on time. Negative
binomial errors accommodate overdispersion. Estimates are summarised as
cumulative relative risks (RR) over the 0--6 month lag window, centred at the
median exposure. Two implementations are provided: a self-contained Python
version and the canonical \texttt{dlnm} R package.

\section{Results}
\subsection{East Java}
Dengue incidence is strongly seasonal, peaking in the wet-season months
(January--March) and lowest in August--September (Figure~\ref{fig:ts}). The
strongest rainfall--dengue cross-correlation occurs at a two-month lag. The
cumulative exposure--response indicates rising dengue risk with increasing
monthly rainfall across the observed range (Figure~\ref{fig:er}). The
negative-binomial DLNM fits the series closely.

\begin{figure}[h]
\centering
\includegraphics[width=0.92\textwidth]{figures/fig1_timeseries.png}
\caption{East Java: monthly dengue cases and climate, 2004--2024
(OpenDengue v1.3).}
\label{fig:ts}
\end{figure}

\begin{figure}[h]
\centering
\includegraphics[width=0.66\textwidth]{figures/fig4_er_rainfall.png}
\caption{East Java cumulative exposure--response for monthly rainfall
(0--6 month lag), relative to the median. Shaded band: 95\% CI.}
\label{fig:er}
\end{figure}

\subsection{Seven provinces}
Extending to a pooled seven-province incidence model preserves the pattern while
revealing between-province heterogeneity. Incidence trajectories differ markedly
in level and timing (Figure~\ref{fig:panel}), yet the peak rainfall--dengue
association clusters at a two-to-three month lag in every province
(Figure~\ref{fig:het}). This consistency of the delayed rainfall effect, against
a backdrop of heterogeneous baseline risk, is exactly the structure an
early-warning system can exploit.

\begin{figure}[h]
\centering
\includegraphics[width=0.95\textwidth]{figures/mp_fig1_incidence_panel.png}
\caption{Monthly dengue incidence per 100{,}000 by province, 2004--2024.}
\label{fig:panel}
\end{figure}

\begin{figure}[h]
\centering
\includegraphics[width=0.7\textwidth]{figures/mp_fig4_heterogeneity.png}
\caption{Between-province heterogeneity: peak rainfall--dengue correlation and
the lag (months) at which it occurs.}
\label{fig:het}
\end{figure}

\subsection{Large-scale drivers: ENSO and IOD}
Local weather is itself modulated by large-scale climate modes. Two were added
to the East Java model as additional cross-bases: the El Ni\~no--Southern
Oscillation, via the Oceanic Ni\~no Index (ONI), and the Indian Ocean Dipole,
via the Dipole Mode Index (DMI) --- the latter especially relevant to Indonesian
rainfall. The estimated exposure--response for ONI (Figure~\ref{fig:oni})
summarises how dengue risk varies across La Ni\~na and El Ni\~no states,
cumulated over a six-month lag. Because these indices partly act \emph{through}
local rainfall, they are naturally collinear with it; whether they add
independent predictive value is an empirical question addressed next.

\begin{figure}[h]
\centering
\includegraphics[width=0.66\textwidth]{figures/fig6_er_oni.png}
\caption{East Java cumulative exposure--response for ENSO (ONI), 0--6 month lag,
relative to neutral conditions. Shaded band: 95\% CI.}
\label{fig:oni}
\end{figure}

\subsection{Forecast skill}
To move from association to a genuine early-warning prototype, the model was
trained on 2004--2019 and evaluated on the held-out 2020--2024 period, against a
seasonal ``climatology'' baseline containing no climate information. For East
Java the climate DLNM improved the out-of-sample forecast substantially
(correlation with observed cases $0.70$ versus $0.38$ for the baseline; an RMSE
skill gain of roughly one fifth) and detected a larger share of true outbreak
months (Figure~\ref{fig:fc}). Adding ENSO and IOD indices at fixed lags did not
improve out-of-sample skill here (correlation $0.65$): in this demonstration the
indices are largely redundant with local rainfall, and entering them instead as
flexible cross-bases \emph{raised} in-sample fit while \emph{lowering} held-out
skill --- a textbook reminder that fit and forecast are different targets. Skill
was also heterogeneous across provinces --- unsurprising given single-point
climate proxies, the COVID-19 disruption to 2020--2021 surveillance, and the
synthetic climate used here. Establishing where and when these drivers add
genuine forecast value, with real reanalysis and seasonal-forecast climate, is
itself a core research question.

\begin{figure}[h]
\centering
\includegraphics[width=0.95\textwidth]{figures/fc_fig1_eastjava_forecast.png}
\caption{East Java out-of-sample forecast. The model is trained only on
2004--2019; the shaded region (2020--2024) is held out. The climate DLNM tracks
the held-out epidemics more closely than the seasonal baseline.}
\label{fig:fc}
\end{figure}

\section{Discussion}
This analysis reproduces, at small scale, the core building block of a dengue
early-warning system: an interpretable, lagged climate--incidence relationship
estimated with a standard model, and shown to be reasonably consistent across
seven provinces. It is deliberately limited: single climate points per province,
no formal forecast evaluation, and surveillance case definitions that vary over
the source period. A doctoral programme would address exactly these gaps ---
fully hierarchical spatio-temporal models, forecast-skill evaluation against
held-out epidemics, incorporation of ENSO, urbanisation and mobility, and
district-level alert thresholds co-designed with health authorities, and
forecast evaluation with proper probabilistic scoring rules and seasonal-forecast
climate rather than the single hold-out shown here.

\section*{Reproducibility}
\small
All code and the real dengue extract are openly available. The analysis runs
end-to-end from the command line; real climate is obtained via a single script
(Open-Meteo, no API key).

\section*{Key references}
\small
Gasparrini A. Distributed lag linear and non-linear models in R: the package
dlnm. \emph{J Stat Softw.} 2011;43(8). \quad
Lowe R, et al. Combined effects of hydrometeorological hazards and urbanisation
on dengue risk in Brazil. \emph{Lancet Planet Health.} 2021. \quad
Clarke J, et al. A global dataset of publicly available dengue case count data.
\emph{Sci Data.} 2024;11:296. \quad
Col\'on-Gonz\'alez FJ, et al. Projecting the risk of mosquito-borne diseases in a
warmer and more populated world. \emph{Lancet Planet Health.} 2021.

\end{document}
